\begin{SCn}
\begin{small}

\scnheader{Protégé}
\scnidtf{Инструмент для разработки OWL-онтологий, разработанный Стэнфордским университетом}
\scnidtf{Редактор онтологий, разработанный в Стэнфорде}
\scniselement{инструмент для разработки онтологий}
\begin{scnindent}
    \scnidtf{редактор онтологий}
\end{scnindent}
\scniselement{семантическое программное обеспечение}
\scnrelfrom{разработчик}{Стэнфордский центр биомедицинской информатики}
\scnrelfrom{год создания}{1987}
\scntext{текущая стабильная версия}{5.6.5}
\scntext{дата выпуска текущей версии}{28 октября 2024}
\scnrelfrom{лицензия}{BSD 2-clause}
\scnrelfrom{язык разработки}{Java}
\scntext{назначение}{Создание и управление сложными семантическими моделями для задач искусственного интеллекта и систем принятия решения}
\begin{scnrelfromset}{варианты исполнения}
    \scnitem{Protégé Desktop}
    \begin{scnindent}
        \scnidtf{Настольный вариант реализации Protégé}
        \scnnote{Распространяется как ZIP-архив с встроенной JRE, требует 1.5 ГБ дискового пространства}
    \end{scnindent}
    \scnitem{WebProtégé}
    \begin{scnindent}
        \scnidtf{Облачный вариант реализации Protégé}
        \scnnote{Реализовано как клиент-серверное решение с REST API и WebSocket для реального взаимодействия. Более 68,000 проектов и 50,000 пользователей, включая команды из MIT и NASA}
    \end{scnindent}
\end{scnrelfromset}
\begin{scnrelfromset}{функциональные возможности}
    \scnfileitem{редактирование онтологий}
    \begin{scnindent}
        \scnidtf{создание и модификация элементов онтологий}
        \scntext{пример использования}{Создание иерархии медицинских диагнозов с указанием отношений "is-a" и "part-of"}
    \end{scnindent}
    \scnfileitem{визуализация онтологий}
    \begin{scnindent}
        \scnidtf{многорежимное отображение структуры онтологий}
        \begin{scnindent}
            \begin{scnrelfromset}{режимы визуализации}
                \scnfileitem{древовидное представление таксонов}
                \scnfileitem{графовые диаграммы отношений}
                \scnfileitem{матрицы свойств}
                \scnfileitem{интерактивные дендрограммы}
            \end{scnrelfromset}
        \end{scnindent}
        \scntext{пример использования}{Визуальная диагностика циклических зависимостей в онтологии транспортных сетей}
    \end{scnindent}
    \scnfileitem{валидация онтологий}
    \begin{scnindent}
        \begin{scnrelfromset}{типы проверок}
            \scnfileitem{синтаксическая корректность OWL}
            \scnfileitem{логическая противоречивость}
            \scnfileitem{нарушения constraints}
            \scnfileitem{соответствие профилям OWL 2 (EL, QL, RL)}
        \end{scnrelfromset}
    \end{scnindent}
    \scnfileitem{поддержка совместной работы}
    \begin{scnindent}
        \begin{scnrelfromset}{уточнение}
            \scnfileitem{отслеживание всех изменений}
            \scnfileitem{комментирование элементов репозитория}
            \scnfileitem{чат с другими пользователями}
            \scnfileitem{голосование по изменениям}
            \scnfileitem{аннотирование классов и экземпляров}
        \end{scnrelfromset}
    \end{scnindent}
\scnfileitem{поддержка плагинов}
\begin{scnindent}
    \scntext{уточнение}{имеет архитектуру, расширяемую за счет плагинов}
    \begin{scnrelfromset}{категории плагинов}
        \scnfileitem{плагины для визуализации (OntoGraf, VOWL)}
        \scnfileitem{плагины для логического вывода (Protege-DL)}
        \scnfileitem{планины для импорта и экспорта (CSV, JSON-LD)}
    \end{scnrelfromset}
    \scnrelfrom{пример}{SHACL4Protege}
    \begin{scnindent}
        \scntext{описание}{плагин для валидации ограничений SHACL}
    \end{scnindent}
\end{scnindent}
\end{scnrelfromset}
\begin{scnrelfromset}{поддерживаемые языки}
    \scnitem{OWL 2}
    \scnitem{OWL}
    \scnitem{RDF}
    \scnitem{RDFS}
    \scnitem{OBO}
    \scnitem{Turtle}
\end{scnrelfromset}
\begin{scnrelfromset}{экосистема расширений}
    \scnitem{OWLViz}
    \begin{scnindent}
        \scntext{назначение}{Визуализация таксономических деревьев с экспортом в векторную графику}
    \end{scnindent}
    \scnitem{SWRLTab}
    \begin{scnindent}
        \scnidtf{Редактор правил Semantic Web Rule Language с синтаксическим анализом}
    \end{scnindent}
    \scnitem{RDFSync}
    \begin{scnindent}
        \scntext{особенность}{Синхронизация с распределенными RDF-хранилищами через SPARQL 1.1}
    \end{scnindent}
\end{scnrelfromset}
\begin{scnrelfromset}{преимущества}
    \scnfileitem{открытый исходный код}
    \begin{scnindent}
        \scntext{уточнение}{Позволяет адаптировать инструмент под специфические задачи, например, интеграцию с медицинскими базами данных через пользовательские плагины.}
    \end{scnindent}
    \scnfileitem{бесплатное использование}
    \scnfileitem{активное сообщество пользователей и разработчиков}
    \scnfileitem{полная поддержка стандартов W3C}
    \scnfileitem{расширяемость через плагины}
    \scnfileitem{удобный интерфейс для начинающих}
\end{scnrelfromset}
\begin{scnrelfromset}{недостатки}
    \scnfileitem{ограниченные возможности интеграции с базами данных без дополнительных плагинов}
    \scnfileitem{относительно ограниченные возможности для разработки корпоративных приложений}
\end{scnrelfromset}
\begin{scnrelfromset}{библиографические источники}
    \scnfileitem{Musen M. A. The Protégé Project: A Look Back and a Look Forward // AI Matters. – 2015. – Vol. 1, N 4. – P. 4–12.}
    \scnfileitem{Noy N. F., Sintek M., Decker S., Crubezy M., Fergerson R. W., Musen M. A. Creating Semantic Web Contents with Protégé-2000 // IEEE Intelligent Systems. – 2001. – Vol. 16, N 2. – P. 60–71.}
    \scnfileitem{Protégé [Электронный ресурс]. — 2025. — URL: https://protege.stanford.edu (дата обращения: 12.05.2025).}
    \scnfileitem{WebProtégé: Collaborative Ontology Development for the Semantic Web [Электронный ресурс]. — 2025. — URL: https://webprotege.stanford.edu (дата обращения: 12.05.2025).}
    \scnfileitem{OWL 2 Web Ontology Language Document Overview [Электронный ресурс]. — 2025. — URL: https://www.w3.org/TR/owl2-overview/ (дата обращения: 12.05.2025).}
\end{scnrelfromset}

\scnheader{Google Knowledge Graph}
\scnidtf{База знаний, используемая Google для улучшения результатов поиска семантически связанной информацией}
\scnidtf{KG (Knowledge Graph)}
\scniselement{формат представления знаний}
\scniselement{семантическая платформа}
\begin{scnrelfromlist}{разработчик}
    \scnitem{Google Inc.}
\end{scnrelfromlist}
\scnrelfrom{год создания}{2012}
\scnrelfrom{основная парадигма}{граф знаний, семантические сети}
\scntext{описание}{Google Knowledge Graph --- технология хранения и представления структурированных знаний в виде графа сущностей и отношений между ними. Каждая сущность описывается набором атрибутов и связей с другими сущностями.}
\scntext{описание}{Knowledge Graph содержит миллиарды фактов о миллионах сущностей: людях, местах, организациях, событиях, продуктах и других объектах реального мира. Данные собираются из различных источников, включая Wikipedia, Wikidata, CIA World Factbook и партнёрские источники.}
\scntext{назначение}{Предназначен для предоставления пользователям прямой информации в результатах поиска, улучшения релевантности поиска через понимание семантики запросов, связывания связанных сущностей и фактов.}
\begin{scnrelfromset}{основные конструкции}
    \scnfileitem{сущность (entity)}
    \scnfileitem{отношение (relation)}
    \scnfileitem{атрибут (attribute)}
    \scnfileitem{триплет (subject-predicate-object)}
\end{scnrelfromset}
\scnrelfrom{модель данных}{ориентированный граф знаний}
\begin{scnrelfromlist}{поддержка}
    \scnfileitem{мультиязычность}
    \scnfileitem{связывание сущностей}
    \scnfileitem{верификация фактов}
    \scnfileitem{масштабируемость}
\end{scnrelfromlist}
\begin{scnrelfromlist}{применение}
    \scnitem{поисковые системы}
    \scnitem{семантический поиск}
    \scnitem{интеллектуальные помощники}
    \scnitem{рекомендательные системы}
    \scnitem{обработка естественного языка}
    \scnitem{извлечение знаний}
\end{scnrelfromlist}
\begin{scnrelfromset}{примеры использования}
    \scnfileitem{Панели знаний (Knowledge Panels) в результатах поиска Google}
    \scnfileitem{Автозаполнение поисковых запросов с учётом семантики}
    \scnfileitem{Связывание биографических данных известных личностей}
    \scnfileitem{Представление информации о фильмах, книгах, музыкальных произведениях}
    \scnfileitem{Геолокационные данные и информация о местах}
    \scnfileitem{Интеграция с Google Assistant для ответов на вопросы}
    \scnfileitem{Улучшение поиска через понимание намерений пользователя}
\end{scnrelfromset}
\begin{scnrelfromset}{технологии и стандарты}
    \scnitem{RDF (Resource Description Framework)}
    \scnitem{Schema.org}
    \scnitem{JSON-LD}
    \scnitem{Microdata}
    \scnitem{Wikidata}
\end{scnrelfromset}
\begin{scnrelfromset}{особенности}
    \scnfileitem{огромный масштаб (миллиарды фактов)}
    \scnfileitem{постоянное обновление и расширение}
    \scnfileitem{мультиязычная поддержка}
    \scnfileitem{интеграция с поисковыми алгоритмами}
    \scnfileitem{использование машинного обучения для извлечения знаний}
    \scnfileitem{поддержка структурированной разметки для веб-мастеров}
\end{scnrelfromset}
\begin{scnrelfromset}{ограничения}
    \scnfileitem{закрытость внутренней структуры и API}
    \scnfileitem{зависимость от качества источников данных}
    \scnfileitem{возможные ошибки в автоматически извлечённых данных}
    \scnfileitem{ограниченный прямой доступ для внешних разработчиков}
\end{scnrelfromset}
\begin{scnrelfromset}{библиографические источники}
    \scnfileitem{Singhal A. Introducing the Knowledge Graph: Things, Not Strings // Official Google Blog. – 2012. – URL: https://blog.google/products/search/introducing-knowledge-graph-things-not/ (дата обращения: 12.05.2025).}
    \scnfileitem{Noy N. F., Gao Y., Jain A., Narayanan A., Patterson A., Taylor J. Industry-scale Knowledge Graphs: Lessons and Challenges // Communications of the ACM. – 2019. – Vol. 62, N 8. – P. 36–43.}
    \scnfileitem{Knowledge Graph Search API [Электронный ресурс] / Google LLC. — 2025. — URL: https://developers.google.com/knowledge-graph (дата обращения: 12.05.2025).}
    \scnfileitem{Schema.org [Электронный ресурс]. — 2025. — URL: https://schema.org (дата обращения: 12.05.2025).}
    \scnfileitem{Wikidata [Электронный ресурс]. — 2025. — URL: https://www.wikidata.org (дата обращения: 12.05.2025).}
\end{scnrelfromset}

\bigskip

\scnheader{Сравнение инструментов представления знаний Protégé и Google Knowledge Graph}
\begin{scneqtovector}
    \scnitem{Protégé}
    \scnitem{Google Knowledge Graph}
\end{scneqtovector}
\begin{scnrelfromset}{сходства}
    \scnfileitem{Оба предназначены для представления и управления структурированными знаниями}
    \scnfileitem{Используют семантические технологии и онтологический подход}
    \scnfileitem{Работают с сущностями и отношениями между ними}
    \scnfileitem{Поддерживают стандарты Semantic Web (RDF, OWL)}
    \scnfileitem{Применяются для организации знаний в машиночитаемом формате}
    \scnfileitem{Позволяют связывать информацию из различных источников}
    \scnfileitem{Используются для улучшения поиска и навигации по данным}
\end{scnrelfromset}
\begin{scnrelfromset}{различия}
    \scnfileitem{Protégé --- инструмент для создания и редактирования онтологий, Google Knowledge Graph --- готовая платформа и база знаний}
    \scnfileitem{Protégé ориентирован на индивидуальных разработчиков и исследователей, Knowledge Graph --- на массовое использование в поисковой системе}
    \scnfileitem{Protégé имеет открытый исходный код и бесплатен, Knowledge Graph --- проприетарная закрытая система Google}
    \scnfileitem{Protégé позволяет создавать собственные онтологии с нуля, Knowledge Graph предоставляет готовую базу знаний}
    \scnfileitem{Protégé работает локально или в веб-версии с ограниченными данными, Knowledge Graph содержит миллиарды фактов о миллионах сущностей}
    \scnfileitem{Knowledge Graph интегрирован с поисковыми алгоритмами Google, Protégé --- независимый инструмент}
    \scnfileitem{Protégé поддерживает ручное создание и редактирование онтологий, Knowledge Graph использует автоматическое извлечение знаний из множества источников}
    \scnfileitem{Knowledge Graph предоставляет публичный API для поиска сущностей, Protégé не имеет аналогичного публичного интерфейса}
\end{scnrelfromset}
\begin{scnrelfromset}{рекомендации по выбору}
    \scnfileitem{Protégé рекомендуется для исследовательских проектов, создания специализированных онтологий и образовательных целей}
    \scnfileitem{Protégé подходит для разработки предметно-ориентированных баз знаний с полным контролем над структурой}
    \scnfileitem{Protégé целесообразно использовать при необходимости экспорта онтологий в стандартных форматах (OWL, RDF)}
    \scnfileitem{Google Knowledge Graph рекомендуется для улучшения видимости контента в поиске Google через структурированную разметку}
    \scnfileitem{Knowledge Graph API полезен для приложений, требующих доступа к обширной базе общих знаний}
    \scnfileitem{Для корпоративных проектов с собственными данными целесообразно создание внутренней базы знаний по аналогии с Knowledge Graph}
    \scnfileitem{Использование Schema.org и JSON-LD рекомендуется для интеграции с Knowledge Graph и другими семантическими платформами}
\end{scnrelfromset}

\end{small}
\end{SCn}
